\documentclass[11pt]{article}
\usepackage[T1]{fontenc}
\usepackage{amsmath,amssymb,amsthm}
\usepackage[margin=1in]{geometry}
\usepackage[hidelinks]{hyperref}\hypersetup{pdftitle={The Zero-Free Columns of the Character Table of S44},pdfauthor={Xiaojian Zeng},pdfsubject={Symmetric-group character tables}}

\newtheorem{theorem}{Theorem}
\newtheorem{lemma}[theorem]{Lemma}
\newcommand{\ZF}{\operatorname{ZF}}

\title{The Zero-Free Columns of the Character Table of $S_{44}$}
\author{Xiaojian Zeng\\
PBC School of Finance, Tsinghua University, 100190\\
\texttt{xiaojianzeng@foxmail.com}}
\date{}

\begin{document}
\maketitle

\begin{abstract}
For a partition $\mu$ of $44$, let the column indexed by $\mu$ in the
ordinary irreducible character table of $S_{44}$ be called zero-free when
$\chi^\lambda(\mu)\ne0$ for every $\lambda\vdash44$.  We determine this
finite layer exactly: the identity column is the only zero-free column, so
$\ZF(44)=\{(1^{44})\}$ and $D(44)=1$.  The proof combines the published
core-distance reduction with two explicit, exactly verified zero witnesses.
The result concerns only $n=44$ and makes no assertion about the asymptotic
zero-free-column conjecture or about priority.
\end{abstract}

\section{Context and statement}

The rows and columns of the ordinary irreducible character table of the
symmetric group $S_n$ are indexed by partitions $\lambda,\mu\vdash n$; the
entry $\chi^\lambda(\mu)$ is the value of the irreducible character afforded
by the Specht module $S^\lambda$ on the class of cycle type $\mu$.  Following
Defant, Hariharan, Lau, and Ono~\cite{DefantEtAl2026}, write
\[
 \ZF(n)=\{\mu\vdash n: \chi^\lambda(\mu)\ne0
          \text{ for every }\lambda\vdash n\},
 \qquad D(n)=|\ZF(n)|.
\]
Mart\'in Duro's question about the integers for which every nonidentity
column contains a zero is the origin of this finite problem
\cite[Question~4.17]{MartinDuro2026}.  Defant et al. give a complete table
only through $n=43$, where $D(43)=7$~\cite[Section~8.2]{DefantEtAl2026}.
The present note supplies the adjacent value $n=44$.

\begin{theorem}\label{thm:main}
For the ordinary irreducible complex character table of $S_{44}$,
\[
             \ZF(44)=\{(1^{44})\},\qquad D(44)=1.
\]
\end{theorem}

The bounded literature comparison above identifies the closest inspected
results; it is not a claim of priority.  In particular, Theorem~\ref{thm:main}
does not address the conjectural asymptotic size of $D(n)$.

\section{Exhaustive reduction}

The reduction of Defant et al.~\cite[Proposition~3.1 and
Corollaries~4.1--4.2]{DefantEtAl2026}, based on the
Murnaghan--Nakayama rule~\cite{Murnaghan1937,Nakayama1940}, says that a
zero-free type for $n\ge3$ must have the form
\[
                    (3^a,2^b,1^c),
\]
where $b$ is even, $b\le s_2(n)$, and $a\le s_3(n)$.  Here $s_2(n)$ is the
least $s\ge0$ for which $n-2s$ is triangular, and $s_3(n)$ is the least
$s\ge0$ for which there is a $3$-core partition of $n-3s$.

For $n=44$, the integers $44,42,40,38$ lie strictly between the consecutive
triangular numbers $36$ and $45$, while
\[
                    44-2\cdot4=36=\frac{8\cdot9}{2}.
\]
Thus $s_2(44)=4$.  The existence criterion for $3$-cores used in the same
paper gives $s_3(44)=0$, since
\[
                3\cdot44+1=133=9^2+9\cdot4+4^2.
\]
It follows that every zero-free column is one of the three types
\begin{equation}\label{eq:candidates}
       (1^{44}),\qquad (2^2,1^{40}),\qquad (2^4,1^{36}).
\end{equation}
This is an exhaustive reduction, not a computationally sampled list.

\section{Exact zero witnesses}

The two nonidentity types in~\eqref{eq:candidates} have the following zero
witnesses:
\begin{center}
\begin{tabular}{c|c|c}
cycle type $\mu$ & row label $\lambda$ & $\chi^\lambda(\mu)$\\ \hline
$(2^2,1^{40})$ & $(14,7,6,5,3,3,2,2,2)$ & $0$\\
$(2^4,1^{36})$ & $(14,11,4,3,3,1,1,1,1,1,1,1,1,1)$ & $0$
\end{tabular}
\end{center}
Both displayed row labels are partitions of $44$.  For completeness, we
describe an exact certificate for the two values.

Under the Frobenius characteristic map, one has
\[
 s_\lambda=\sum_{\nu\vdash44}\frac{\chi^\lambda(\nu)}{z_\nu}p_\nu,
\]
and the Jacobi--Trudi identity gives
$s_\lambda=\det(h_{\lambda_i-i+j})$; see
Macdonald~\cite[Chapter~I]{Macdonald1995}.  Specialize $p_1=1$, $p_2=y$,
and $p_k=0$ for $k>2$.  Then
\[
 \sum_{m\ge0}h_m t^m=\exp(t+yt^2/2),\qquad
 h_m=\sum_{r=0}^{\lfloor m/2\rfloor}
       \frac{y^r}{(m-2r)!2^r r!}.
\]
Consequently, if $\mu_b=(2^b,1^{44-2b})$, then
\begin{equation}\label{eq:coefficient}
 \chi^\lambda(\mu_b)
 =(44-2b)!2^b b!\,[y^b]\det(h_{\lambda_i-i+j}).
\end{equation}

The determinant in~\eqref{eq:coefficient} was evaluated in the truncated
ring $\mathbb Q[y]/(y^{b+1})$.  A directly reproducible row-subset recurrence
starts with $F_0(\varnothing)=1$ and, when column $j$ is appended after the
set $S$ of previously used columns, adds
\[
 (-1)^{|\{s\in S:s>j\}|}F_i(S)h_{\lambda_{i+1}-(i+1)+(j+1)}
 \quad\text{to}\quad F_{i+1}(S\cup\{j\}).
\]
All polynomial coefficients are rational numbers.  For the first displayed
partition with $b=2$, and for the second with $b=4$, the coefficient
$[y^b]F_\ell(\{0,\ldots,\ell-1\})$ is exactly $0$.  Equation
\eqref{eq:coefficient} therefore proves both character values in the table,
without floating-point evaluation.

Finally, the identity type survives universally:
\[
              \chi^\lambda(1^{44})=\dim S^\lambda>0
              \qquad(\lambda\vdash44).
\]
Thus the only surviving candidate in~\eqref{eq:candidates} is $(1^{44})$,
which proves Theorem~\ref{thm:main}.

\section{Computational verification and reproducibility}

The witnesses were discovered by a second exact implementation of the
Murnaghan--Nakayama recurrence.  It generated all $75{,}175$ partitions of
$44$ and evaluated both reduced nonidentity columns using beta-number rim-hook
moves and the integral hook-length formula at the remaining fixed points.  It
found $18$ zero rows for $(2^2,1^{40})$ and $2$ for $(2^4,1^{36})$.  The
Jacobi--Trudi computation above uses neither beta numbers nor rim-hook
enumeration and independently verifies the only two values needed in the
proof.  A third regression recomputed every published value $D(n)$ for
$1\le n\le43$ and the seven published survivors at $n=43$.

This separation between theorem-bearing exact identities and diagnostic
finite checks is useful in other partition-based computations as well.  For
example, Zeng~\cite[Section~4]{Zeng2026MultiplicityArxiv} describes an exact
integer cut-partition dynamic program checked by direct small-order
enumeration, while explicitly keeping those computations outside the premise
of the theorem.  Here likewise, the historical regression supports the
implementation but is not needed once the two exact coefficients in
\eqref{eq:coefficient} have been verified.

The accompanying files \texttt{mn\_scan.py},
\texttt{independent\_verify.py}, and \texttt{reproduce\_table.py}, together
with their plain-text transcripts, contain the exact computations described
above.  They require only Python's standard library.

\section*{Computational and AI assistance disclosure}

The author used OpenAI Codex to assist with the finite search, exact-checking
workflow, literature screening, and manuscript preparation.  The decisive
character values were computed with exact integer or rational arithmetic and
were checked by structurally different Murnaghan--Nakayama and
Jacobi--Trudi implementations.  No floating-point calculation or AI output is
used as a mathematical premise.

\bibliographystyle{alpha}
\bibliography{references}

\end{document}
